% =============================================================================
%  PARTIE I — Guide utilisateur
% =============================================================================

\chapter{Présentation}
\label{ch:presentation}

\Icare{} est un petit contrôleur de lumière DMX qui se présente comme une
application web locale. Concrètement : un programme Python démarre un serveur
\emph{Flask} sur votre machine, et votre navigateur affiche l'interface. Un
adaptateur USB-DMX branché en USB assure la liaison physique avec les
projecteurs.

Le projet vise les \textbf{groupes amateurs} qui veulent se passer d'un vrai
pupitre et d'un éclairagiste dédié : on branche, on ouvre la page, et les
lumières répondent.

\section{Principes directeurs}

\begin{itemize}[leftmargin=1.4em]
  \item \textbf{100 \% local.} Rien ne part sur Internet ; le navigateur ne
        parle qu'à \code{127.0.0.1:5000}. Pas de compte, pas de cloud.
  \item \textbf{Une seule interface.} Le navigateur est le pupitre : trois
        onglets suffisent (\emph{Direct}, \emph{Séquence}, \emph{Live}).
  \item \textbf{Démarrage rapide.} Une commande pour installer, une commande pour
        lancer. Une version \code{.exe} autonome existe aussi (voir
        chapitre~\ref{ch:installation}).
  \item \textbf{Le spectaculaire accessible.} Effets intégrés (arc-en-ciel, feu,
        gyrophare, chenillard), mode \emph{audio-réactif} au microphone, et
        lecture calée au tempo avec le bouton \emph{TAP}.
\end{itemize}

\section{Vocabulaire}

Avant d'aller plus loin, fixons les mots employés dans ce document et dans
l'interface.

\begin{description}[leftmargin=1.6em,style=nextline]
  \item[Univers (DMX)] Ensemble des 512 canaux d'une ligne DMX. Chaque canal
        porte une valeur de 0 à 255.
  \item[Module / projecteur (\emph{fixture})] Un appareil branché sur la ligne,
        qui occupe un ou plusieurs canaux. Dans \Icare{}, les modules portent des
        noms comme \code{ParLED1} ou \code{Strobe1}.
  \item[Canal de base] Premier canal d'un projecteur. Si un PAR RGB démarre au
        canal 9, alors R, G et B occupent respectivement les canaux 9, 10 et 11.
  \item[Séquence] Une suite d'étapes (\emph{lignes}) décrivant l'état des
        projecteurs, enregistrée dans un fichier JSON.
  \item[Section] Étiquette donnée à un groupe de lignes (\code{Intro\_x4},
        \code{Couplet1}\dots) pour se repérer.
  \item[Beat / ligne] Une unité de temps de la séquence. Une ligne = un temps
        musical.
  \item[BPM] \emph{Beats per minute}. Cadence de lecture de la séquence Live.
\end{description}

\section{Les trois onglets}

La barre de navigation en haut de page donne accès aux trois espaces de travail,
résumés dans le tableau~\ref{tab:onglets}.

\begin{table}[H]
  \centering
  \begin{tabular}{@{}p{2.4cm}p{11.5cm}@{}}
    \toprule
    \textbf{Onglet} & \textbf{À quoi il sert} \\
    \midrule
    \textbf{Direct} & Pilotage immédiat : couleurs, intensité, pan/tilt,
      blackout, couleurs aléatoires, strobe, effets intégrés, audio-réactif et
      visualisation 3D. \\
    \addlinespace
    \textbf{Séquence} & Édition d'un show temps par temps : table de lignes par
      module, insertion d'effets \emph{Uniform}, \emph{Fondu} et \emph{Strobe},
      sauvegarde dans un fichier JSON. \\
    \addlinespace
    \textbf{Live} & Lecture d'une séquence sur scène : choix des modules, BPM
      réglable ou mesuré au \emph{TAP}, transport, barre de progression et vue 3D
      du faisceau. \\
    \bottomrule
  \end{tabular}
  \caption{Rôle des trois onglets.}
  \label{tab:onglets}
\end{table}

En complément, la \textbf{barre supérieure} (ou \emph{topbar}) affiche un
témoin d'état DMX et un bouton de bascule de thème clair/sombre.

\begin{figure}[H]
  \centering
  \begin{tikzpicture}[
      font=\small,
      box/.style={draw=icodeln,fill=icodebg,rounded corners,inner sep=8pt,align=center,minimum width=3.4cm},
      nav/.style={draw=icareblue,fill=icareblue!10,rounded corners,inner sep=6pt,align=center,minimum width=2.2cm},
      arrow/.style={-Latex,draw=icodecmt},
    ]
    \node[box,minimum width=13cm] (top) {Topbar : marque \textbf{Icare} \quad|\quad Accueil \quad Direct \quad Séquence \quad Live \quad|\quad témoin DMX \quad thème};
    \node[nav,below=1.1cm of top,yshift=0.1cm,xshift=-4.2cm] (d) {Direct};
    \node[nav,right=0.5cm of d] (s) {Séquence};
    \node[nav,right=0.5cm of s] (l) {Live};
    \node[box,below=0.9cm of s,minimum width=13cm] (content) {Contenu de l'onglet actif};
    \draw[arrow] (d) -- (content);
    \draw[arrow] (s) -- (content);
    \draw[arrow] (l) -- (content);
  \end{tikzpicture}
  \caption{Organisation générale de l'interface.}
  \label{fig:ui}
\end{figure}

\chapter{Installation et démarrage}
\label{ch:installation}

\section{Prérequis}

\begin{itemize}[leftmargin=1.4em]
  \item Windows (les autres systèmes fonctionnent en développement, mais le
        paquetage fourni cible Windows).
  \item Python 3.10 ou plus récent.
  \item Un adaptateur USB-DMX compatible \emph{Enttec USB Pro} (ou un clone qui
        implémente le même protocole).
\end{itemize}

\section{Démarrage depuis les sources}

\begin{lstlisting}[language=bash,caption={Installation et lancement en développement.}]
python -m venv .venv
.\.venv\Scripts\Activate.ps1
pip install -r requirements.txt
python run_webapp.py
\end{lstlisting}

Le serveur écoute sur \code{http://127.0.0.1:5000}. Ouvrez cette adresse dans
votre navigateur. En mode \emph{exécutable} empaqueté, le navigateur s'ouvre
automatiquement au démarrage.

\section{Connecter l'adaptateur DMX}

\begin{enumerate}[leftmargin=1.6em]
  \item Branchez l'adaptateur USB-DMX sur le PC et reliez-le à la ligne DMX des
        projecteurs (XLR 5 broches ou 3 broches selon votre câblage).
  \item Dans l'onglet \textbf{Direct}, renseignez le \emph{canal de base} du
        premier projecteur (9 par défaut), puis cliquez sur
        \textbf{Connexion DMX}.
  \item Le témoin en haut à droite passe au vert (\emph{DMX OK}).
\end{enumerate}

Si aucun port n'est indiqué, \Icare{} choisit automatiquement le premier port
COM dont la description contient \og USB \fg{}. Vous pouvez forcer un port en
le saisissant dans le champ \emph{Port} (par exemple \code{COM3}).

\section{Construire un exécutable Windows}

Pour distribuer \Icare{} sur un PC sans Python, utilisez le script fourni :

\begin{lstlisting}[language=bash,caption={Construction de l'exécutable autonome.}]
.\build_exe.ps1
\end{lstlisting}

Cela produit \file{dist\textbackslash IcareWeb.exe} (environ 30 Mo). Au
lancement, l'exécutable démarre le serveur, ouvre le navigateur et crée ses
fichiers de travail à côté de lui.

\section{Où \Icare{} écrit-il ses fichiers ?}

\begin{table}[H]
  \centering
  \begin{tabular}{@{}llp{7.4cm}@{}}
    \toprule
    \textbf{Fichier / dossier} & \textbf{Type} & \textbf{Contenu} \\
    \midrule
    \file{Sequences/} & dossiers & Shows en JSON (fournis + créés). \\
    \file{Palettes/} & dossier & Palettes de couleurs nommées. \\
    \file{icare\_settings.json} & fichier & Préférences persistantes (thème, port,
      canal de base, dernière séquence, favoris). \\
    \file{icare.log} & fichier & Journal rotatif (512 Ko $\times$ 3). \\
    \bottomrule
  \end{tabular}
  \caption{Fichiers produits à l'exécution.}
  \label{tab:fichiers}
\end{table}

En développement, ces fichiers apparaissent à la racine du projet. En version
empaquetée, ils sont écrits à côté de l'exécutable, jamais dans l'archive
lecture seule.

\chapter{L'onglet Direct}
\label{ch:direct}

L'onglet \textbf{Direct} est le pupitre manuel : tout ce qu'on y fait est envoyé
immédiatement au contrôleur DMX.

\section{Connexion}

Le panneau \emph{Connexion} contient :

\begin{itemize}[leftmargin=1.4em]
  \item \textbf{Canal de base (R)} : premier canal couleur du projecteur
        (0 à 509).
  \item \textbf{Port (optionnel)} : nom du port COM ; vide = auto-détection.
  \item \textbf{Connexion DMX} : ouvre le port série.
  \item \textbf{Mettre à jour canal de base} : change le canal de base sans
        rouvrir le port.
\end{itemize}

\section{Couleur et intensité}

Les curseurs \textbf{Rouge}, \textbf{Vert} et \textbf{Bleu} vont de 0 à 255 ;
le curseur \textbf{Amplitude} règle la luminosité globale. Le sélecteur de
couleur (color picker) pilote directement les trois canaux RGB. Un aperçu
rectangulaire montre la couleur résultante.

Le bouton \textbf{+ Favori} enregistre la couleur courante dans la palette de
favoris ; un clic sur une pastille la rappelle. La croix \code{x} retire un
favori. La liste est persistée dans les réglages (32 favoris maximum).

\section{Orientation}

Les projecteurs du patch sont \textbf{fixes} : ils n'ont ni pan ni tilt (aucune
tête mobile). Leur direction se règle une fois, au montage, par le champ
\code{angle} de chaque appareil dans le patch. L'onglet Direct ne propose donc
plus de réglage de position ni d'orientation.

\section{Les modes}

\paragraph{Blackout / Stop.} \textbf{Blackout} éteint tout immédiatement
(couleurs à 0, amplitude à 0). \textbf{Stop} arrête le mode en cours.

\paragraph{Couleurs aléatoires.} Trois paramètres :
\emph{fréquence} (Hz), \emph{lissage} (0--100 \%) et \emph{luminosité}
(0--255). Le lissage interpole progressivement vers la couleur cible pour un
rendu plus doux.

\paragraph{Strobe (overlay).} Un strobe qui \emph{se superpose} au mode courant
sans l'interrompre. Réglages : fréquence (Hz), rapport cyclique (\%),
canal DMX optionnel et durée (s). Si un canal est précisé, c'est ce canal qui
clignote ; sinon l'image entière est alternée entre la couleur courante et le
noir.

\paragraph{Bibliothèque d'effets.} Quatre effets procéduraux :
\textbf{Rainbow}, \textbf{Fire}, \textbf{Police} et \textbf{Chase RGB}, avec
réglage de \emph{vitesse} et de \emph{luminosité}. Voir le
chapitre~\ref{ch:effets} pour le détail.

\paragraph{Audio-réactif (bêta).} Le microphone du navigateur pilote
l'intensité au rythme de la musique. On règle la \emph{sensibilité} et la
\emph{couleur} ; une barre indique le niveau détecté. Voir le
chapitre~\ref{ch:audio}.

\section{Visualisation 3D}

Le panneau \emph{Visualisation 3D} dessine un projecteur, son faisceau et un
faisceau spot sur une scène. Couleur, intensité et orientation suivent les
curseurs (et les modes actifs). Pendant un strobe overlay, l'opacité du faisceau
clignote au même rythme.

\chapter{L'onglet Séquence}
\label{ch:sequence-ui}

L'onglet \textbf{Séquence} est l'éditeur de spectacle. Il manipule un fichier
JSON que l'on pourra lire ensuite dans l'onglet \emph{Live}.

\section{Fichiers}

Le panneau \emph{Fichiers} liste les séquences disponibles. Chaque entrée offre
un bouton \textbf{Charger} et une croix \code{x} de suppression. Le champ du
haut permet de nommer puis \textbf{Sauvegarder} le show courant.

\section{Barre d'édition}

\begin{table}[H]
  \centering
  \begin{tabular}{@{}lp{10.5cm}@{}}
    \toprule
    \textbf{Champ} & \textbf{Signification} \\
    \midrule
    \textbf{Module} & Projecteur ciblé par l'insertion d'effet. \\
    \textbf{\#Ligne} & Numéro de la ligne (temps) où insérer. \\
    \textbf{\#Tps} & Nombre de temps que dure l'effet (\code{n}). \\
    \textbf{Fréquence} & Fréquence du strobe, en Hz. \\
    Couleur 1 / Couleur 2 & Couleurs de départ et d'arrivée. \\
    \bottomrule
  \end{tabular}
  \caption{Contrôles de la barre d'édition.}
\end{table}

Trois boutons insèrent des effets :

\begin{itemize}[leftmargin=1.4em]
  \item \textbf{Uniform} : couleur fixe pendant \code{n} temps.
  \item \textbf{Fondu} : dégradé linéaire de la couleur 1 vers la couleur 2 sur
        \code{n} temps.
  \item \textbf{Strobe} : clignotement à la fréquence donnée pendant \code{n}
        temps.
\end{itemize}

Le bouton \textbf{+ Ligne} ajoute une ligne vierge, et \textbf{Effacer cellule}
vide la cellule ciblée.

\section{Édition directe du tableau}

Le tableau comporte une colonne par module. Les cellules de fonction sont
éditionnables directement : on peut y taper à la main un jeton du langage
décrit au chapitre~\ref{ch:dsl}. Le caractère \token{|} prolonge l'effet
commencé sur la ligne précédente.

\begin{note}
Une cellule vide signifie \og{}ne rien changer pour ce module à ce temps\fg{}.
C'est utile pour ne modifier qu'un seul projecteur à la fois.
\end{note}

\chapter{L'onglet Live}
\label{ch:live-ui}

L'onglet \textbf{Live} joue une séquence en conditions de spectacle.

\section{Charger et sélectionner}

\begin{enumerate}[leftmargin=1.6em]
  \item Choisissez une séquence dans la liste (\textbf{Charger}).
  \item Cochez les \textbf{modules} à piloter. Par défaut, tous les modules de la
        séquence sont pris en compte. Les boutons \emph{Tout sélectionner} et
        \emph{Tout désélectionner} font le tri.
\end{enumerate}

\section{Régler le tempo}

Deux méthodes :

\begin{itemize}[leftmargin=1.4em]
  \item Saisir un \textbf{BPM} (20 à 300) dans le champ numérique.
  \item Cliquer sur \textbf{TAP} en rythme : après deux appuis rapprochés,
        \Icare{} calcule le tempo moyen et l'applique. \textbf{Reset} efface
        l'historique des appuis.
\end{itemize}

\section{Transport}

\textbf{Play}, \textbf{Pause}, \textbf{Stop}, \textbf{Précédent} et
\textbf{Suivant} contrôlent la lecture. La barre de
progression, le nom de la section courante et le compteur \code{i / total}
donnent l'avancement. À la fin, la dernière image est maintenue puis la lecture
s'arrête.

\section{Vue 3D}

Le faisceau affiché prend la couleur envoyée au DMX en temps réel, ce qui permet
de vérifier ce qui part sur la ligne même sans regarder les projecteurs.

\chapter{Réglages et persistance}
\label{ch:reglages}

Peu de choses à régler, mais elles survivent aux redémarrages.

\begin{table}[H]
  \centering
  \begin{tabular}{@{}ll@{}}
    \toprule
    \textbf{Réglage} & \textbf{Effet} \\
    \midrule
    Thème & Clair ou sombre (bouton lune/soleil de la topbar). \\
    Dernier port COM & Pré-remplit le champ port. \\
    Canal de base & Restauré au démarrage. \\
    Dernière séquence & Mémorisée pour rechargement. \\
    Couleurs favorites & Palette affichée dans l'onglet Direct. \\
    \bottomrule
  \end{tabular}
  \caption{Réglages persistés.}
\end{table}

Tout est écrit dans \file{icare\_settings.json}. Le journal \file{icare.log}
enregistre les événements (connexion, erreurs série, etc.) et est automatiquement
limité en taille.

\begin{avertissement}
Le fichier de réglages est du JSON : s'il devient illisible (édition manuelle
ratée), \Icare{} le réinitialise silencieusement au démarrage suivant.
\end{avertissement}
